\chapter{2004年上海卷（春）}
\section{填空题}
\begin{problem}
若复数 \(z\) 满足 \(z(1+\i)=2\)，则 \(z\) 的实部是\fillinblank{}.
\end{problem}

\begin{answer}
1
\end{answer}

\begin{solution}
由题意，\(z=\dfrac{2}{1+\i}=\dfrac{2(1-\i)}{(1+\i)(1-\i)}=1-\i\)，所以 \(z\) 的实部为 \(1\).
\end{solution}

\begin{problem}
方程 \(\lg x+\lg(x+3)=1\) 的解 \(x=\)\fillinblank{}.
\end{problem}

\begin{answer}
2
\end{answer}

\begin{solution}
由对数的定义域可知 \(x>0\). 因为 \(\lg x+\lg(x+3)=\lg\left(x(x+3)\right)\)，所以 \(x(x+3)=10\)，即 \(x^2+3x-10=0\). 因此 \((x+5)(x-2)=0\)，结合 \(x>0\)，得 \(x=2\).
\end{solution}

\begin{problem}
在 \(\triangle ABC\) 中，\(a\)，\(b\)，\(c\) 分别是 \(\angle A\)，\(\angle B\)，\(\angle C\) 所对的边. 若 \(\angle A=105^{\circ}\)，\(\angle B=45^{\circ}\)，\(b=2\sqrt{2}\)，则 \(c=\)\fillinblank{}.
\end{problem}

\begin{answer}
2
\end{answer}

\begin{solution}
由三角形内角和，\(\angle C=180^{\circ}-105^{\circ}-45^{\circ}=30^{\circ}\). 根据正弦定理，\(\dfrac{c}{\sin30^{\circ}}=\dfrac{b}{\sin45^{\circ}}\)，所以 \(c=2\sqrt{2}\cdot\dfrac{1/2}{\sqrt{2}/2}=2\).
\end{solution}

\begin{problem}
过抛物线 \(y^{2}=4x\) 的焦点 \(F\) 作垂直于 \(x\) 轴的直线，交抛物线于 \(A\)，\(B\) 两点，则以 \(F\) 为圆心，\(AB\) 为直径的圆方程是\fillinblank{}.
\end{problem}

\begin{answer}
\((x-1)^2+y^2=4\)
\end{answer}

\begin{solution}
抛物线 \(y^2=4x\) 的焦点为 \(F(1,0)\). 过 \(F\) 作垂直于 \(x\) 轴的直线，其方程为 \(x=1\). 代入抛物线方程得 \(y^2=4\)，所以 \(A(1,2)\)，\(B(1,-2)\)，从而 \(AB=4\). 以 \(F\) 为圆心、\(AB\) 为直径的圆半径为 \(2\)，故圆方程为 \((x-1)^2+y^2=4\).
\end{solution}

\begin{problem}
已知函数 \(f(x)=\log_3\left(\frac{4}{x}+2\right)\)，则方程 \(f^{-1}(x)=4\) 的解 \(x=\)\fillinblank{}.
\end{problem}

\begin{answer}
1
\end{answer}

\begin{solution}
由反函数的定义，方程 \(f^{-1}(x)=4\) 等价于 \(x=f(4)\). 因为 \(f(4)=\log_3\left(\dfrac{4}{4}+2\right)=\log_3 3=1\)，所以 \(x=1\).
\end{solution}

\begin{problem}
如图，在底面边长为 \(2\) 的正三棱锥 \(V-ABC\) 中，\(E\) 是 \(BC\) 的中点，若 \(\triangle VAE\) 的面积是 \(\frac{1}{4}\)，则侧棱 \(VA\) 与底面所成角的大小为\fillinblank{}. （结果用反三角函数值表示）
\bitmapfigure[max width=0.32\linewidth]{img/2004/shanghai_spring/q6_fig1.png}
\end{problem}

\begin{answer}
\(\arctan\dfrac{1}{4}\)
\end{answer}

\begin{solution}
设正三角形 \(ABC\) 的中心为 \(O\). 因为 \(V-ABC\) 是正三棱锥，所以 \(VO\perp\) 底面，且 \(O\) 在中线 \(AE\) 上. 在边长为 \(2\) 的正三角形 \(ABC\) 中，\(AE=\sqrt{3}\)，\(AO=\dfrac{2}{3}AE=\dfrac{2\sqrt{3}}{3}\). 由于 \(VO\perp AE\)，有 \(\dfrac{1}{2}AE\cdot VO=\dfrac{1}{4}\)，所以 \(VO=\dfrac{\sqrt{3}}{6}\). 设 \(VA\) 与底面所成的角为 \(\theta\)，则 \(\theta=\angle VAO\)，从而 \(\tan\theta=\dfrac{VO}{AO}=\dfrac{\sqrt{3}/6}{2\sqrt{3}/3}=\dfrac{1}{4}\). 因此 \(\theta=\arctan\dfrac{1}{4}\).
\end{solution}

\begin{problem}
在数列 \(\{a_{n}\}\) 中，\(a_{1}=3\)，且对任意大于 \(1\) 的正整数 \(n\)，点 \((\sqrt{a_n},\sqrt{a_{n-1}})\) 在直线 \(x-y-\sqrt{3}=0\) 上，则 \(\displaystyle\lim_{n\to\infty} \frac{a_n}{(n+1)^2}=\)\fillinblank{}.
\end{problem}

\begin{answer}
3
\end{answer}

\begin{solution}
由点 \((\sqrt{a_n},\sqrt{a_{n-1}})\) 在直线 \(x-y-\sqrt{3}=0\) 上，得 \(\sqrt{a_n}-\sqrt{a_{n-1}}=\sqrt{3}\). 又因为 \(a_1=3\)，所以 \(\sqrt{a_1}=\sqrt{3}\). 逐项相加可得 \(\sqrt{a_n}=n\sqrt{3}\)，即 \(a_n=3n^2\). 因而 \(\displaystyle\lim_{n\to\infty}\frac{a_n}{(n+1)^2}=\lim_{n\to\infty}3\left(\frac{n}{n+1}\right)^2=3\).
\end{solution}

\begin{problem}
根据下列 \(5\) 个图形及相应点的个数的变化规律，试猜测第 \(n\) 个图中有\fillinblank{}个点.
\FigureLayoutDeclare{img/2004/shanghai_spring/q8_fig1.png}{12.0cm}{10bp 9bp 9bp 10bp}
\bitmapfigure[max width=0.76\linewidth]{img/2004/shanghai_spring/q8_fig1.png}
\end{problem}

\begin{answer}
\(n^2-n+1\)
\end{answer}

\begin{solution}
设第 \(n\) 个图形中的点数为 \(u_n\). 由图可知，第一个图形有 \(u_1=1\) 个点；从第 \(k\) 个图形增加到第 \(k+1\) 个图形时，在原图形的两侧各增加 \(k\) 个点，所以增加 \(2k\) 个点. 因此 \(u_{k+1}=u_k+2k\)，从而
\(u_n=1+2\left(1+2+\cdots+(n-1)\right)=1+n(n-1)=n^2-n+1\).
\end{solution}

\begin{problem}
一次二期课改经验交流会打算交流试点学校的论文 \(5\) 篇和非试点学校的论文 \(3\) 篇. 若任意排列交流次序，则最先和最后交流的论文都为试点学校的概率是\fillinblank{}. （结果用分数表示）
\end{problem}

\begin{answer}
\(\dfrac{5}{14}\)
\end{answer}

\begin{solution}
任意排列时，最先交流的论文为试点学校论文的概率是 \(\dfrac{5}{8}\). 在最先交流的论文已经是试点学校论文的条件下，剩下 \(7\) 篇论文中还有 \(4\) 篇试点学校论文，因此最后交流的论文也是试点学校论文的条件概率为 \(\dfrac{4}{7}\). 所求概率为 \(\dfrac{5}{8}\cdot\dfrac{4}{7}=\dfrac{5}{14}\).
\end{solution}

\begin{problem}
若平移椭圆 \(4(x+3)^2+9y^2=36\)，使平移后的椭圆中心在第一象限，且它与 \(x\) 轴，\(y\) 轴分别只有一个交点，则平移后的椭圆方程是\fillinblank{}.
\end{problem}

\begin{answer}
\(\dfrac{(x-3)^2}{9}+\dfrac{(y-2)^2}{4}=1\)
\end{answer}

\begin{solution}
设平移后的椭圆中心为 \((h,k)\)，则 \(h>0\)，\(k>0\)，且椭圆方程为 \(\dfrac{(x-h)^2}{9}+\dfrac{(y-k)^2}{4}=1\). 与 \(x\) 轴只有一个交点，意味着令 \(y=0\) 后所得方程只有一个解，因此椭圆与 \(x\) 轴相切. 代入得 \(\dfrac{(x-h)^2}{9}+\dfrac{k^2}{4}=1\)，要只有一个解必须有 \(k=2\). 同理，与 \(y\) 轴只有一个交点意味着令 \(x=0\) 后相切，故 \(h=3\). 所以平移后的椭圆方程为 \(\dfrac{(x-3)^2}{9}+\dfrac{(y-2)^2}{4}=1\).
\end{solution}

\begin{problem}
如图，在由二项式系数所构成的杨辉三角形中，第\fillinblank{}行中从左至右第\(14\)与第\(15\)个数的比为\(2:3\).

\begin{center}
\small
\begin{tabular}{@{}lcccccc@{}}
第\(0\)行 & \(1\) & & & & & \\
第\(1\)行 & \(1\) & \(1\) & & & & \\
第\(2\)行 & \(1\) & \(2\) & \(1\) & & & \\
第\(3\)行 & \(1\) & \(3\) & \(3\) & \(1\) & & \\
第\(4\)行 & \(1\) & \(4\) & \(6\) & \(4\) & \(1\) & \\
第\(5\)行 & \(1\) & \(5\) & \(10\) & \(10\) & \(5\) & \(1\) \\
\(\cdots\) & & & & & &
\end{tabular}
\end{center}
\end{problem}

\begin{answer}
34
\end{answer}

\begin{solution}
按题中从第 \(0\) 行开始的编号，第 \(n\) 行的第 \(14\)、第 \(15\) 个数分别为 \(\mathrm C_n^{13}\)、\(\mathrm C_n^{14}\). 因而
\(\dfrac{\mathrm C_n^{13}}{\mathrm C_n^{14}}=\dfrac{14}{n-13}=\dfrac{2}{3}\)，解得 \(n-13=21\)，所以 \(n=34\).
\end{solution}

\begin{problem}
在等差数列 \(\{a_{n}\}\) 中，当 \(a_{r}=a_{s}\)（\(r\neq s\)）时，\(\{a_{n}\}\) 必定是常数数列. 然而在等比数列 \(\{a_{n}\}\) 中，对某些正整数 \(r\)，\(s\)（\(r\neq s\)），当 \(a_{r}=a_{s}\) 时，非常数数列 \(\{a_{n}\}\) 的一个例子是\fillinblank{}.
\end{problem}

\begin{answer}
\(a,-a,a,-a,\ldots\)（\(a\ne0\)，\(r\) 与 \(s\) 同为奇数或偶数）
\end{answer}

\begin{solution}
取首项为 \(a\)，公比为 \(-1\)，其中 \(a\ne0\). 则所得等比数列为 \(a,-a,a,-a,\ldots\)，它不是常数数列. 对于两个不同且同奇偶性的正整数 \(r\)、\(s\)，有 \((-1)^{r-1}=(-1)^{s-1}\)，所以 \(a_r=a_s\). 这就是所求例子.
\end{solution}

\section{单选题}
\begin{problem}
下列函数中，周期为 \(1\) 的奇函数是（\quad）

\choices
    {\(y = 1 - 2\sin^2\pi x\)}
    {\( y = \sin \left(2\pi x + \frac{\pi}{3}\right) \)}
    {\(y = \tan \frac{\pi}{2} x\)}
    {\( y = \sin \pi x \cos \pi x \)}
\end{problem}

\begin{answer}
D
\end{answer}

\begin{solution}
选项 A 中，\(1-2\sin^2\pi x=\cos2\pi x\) 是偶函数；选项 B 虽然周期为 \(1\)，但有相位平移，不是奇函数；选项 C 中 \(\tan\dfrac{\pi}{2}x\) 的周期为 \(2\)，不是周期为 \(1\)；选项 D 中 \(\sin\pi x\cos\pi x=\dfrac{1}{2}\sin2\pi x\)，既是奇函数，周期又为 \(1\). 所以选 D.
\end{solution}

\begin{problem}
若非空集合 \(M \subseteq N\)，则“\(a \in M\) 或 \(a \in \N^{*}\)”是“\(a \in M \cap N^{*}\)”的（\quad）

\choices
    {充分非必要条件}
    {必要非充分条件}
    {充要条件}
    {既非充分又非必要条件}
\end{problem}

\begin{answer}
B
\end{answer}

\begin{solution}
若 \(a\in M\cap\N^{*}\)，则同时有 \(a\in M\) 和 \(a\in\N^{*}\)，所以命题“\(a\in M\) 或 \(a\in\N^{*}\)”必然成立. 反过来，前一个命题成立时，可能只有 \(a\in M\)，也可能只有 \(a\in\N^{*}\)，不能保证 \(a\) 同时属于两者的交集. 因此前一个命题是后一个命题的必要非充分条件，选 B.
\end{solution}

\begin{problem}
在 \(\triangle ABC\) 中，有命题

\begin{circlelist}
\item \(\overrightarrow{AB}-\overrightarrow{AC}=\overrightarrow{BC}\)；
\item \(\overrightarrow{AB}+\overrightarrow{BC}+\overrightarrow{CA}=\overrightarrow{0}\)；
\item 若 \((\overrightarrow{AB}+\overrightarrow{AC})\cdot(\overrightarrow{AB}-\overrightarrow{AC})=0\)，则 \(\triangle ABC\) 为等腰三角形；
\item 若 \(\overrightarrow{AC}\cdot\overrightarrow{AB}>0\)，则 \(\triangle ABC\) 为锐角三角形.
\end{circlelist}

上述命题正确的是（\quad）

\choices
    {\circled{1}\circled{2}}
    {\circled{1}\circled{4}}
    {\circled{2}\circled{3}}
    {\circled{2}\circled{3}\circled{4}}
\end{problem}

\begin{answer}
C
\end{answer}

\begin{solution}
由向量减法，\(\overrightarrow{AB}-\overrightarrow{AC}=\overrightarrow{CB}\)，所以命题\(\circled{1}\)错误. 命题\(\circled{2}\)是三角形的向量首尾相接关系，因而正确. 对命题\(\circled{3}\)，有
\((\overrightarrow{AB}+\overrightarrow{AC})\cdot(\overrightarrow{AB}-\overrightarrow{AC})=|\overrightarrow{AB}|^2-|\overrightarrow{AC}|^2\). 该式等于零时，\(AB=AC\)，所以三角形为等腰三角形，命题\(\circled{3}\)正确. 命题\(\circled{4}\)只说明 \(\angle A\) 为锐角，不能推出另外两个角也为锐角，所以命题\(\circled{4}\)错误. 因此正确命题为\(\circled{2}\)、\(\circled{3}\)，选 C.
\end{solution}

\begin{problem}
若 \(p = a + \frac{1}{a} + 2\)（\(a>0\)），\(q = \arccos t\)（\(-1 \leqslant t \leqslant 1\)），则下列不等式恒成立的是（\quad）

\choices
    {\(p \geqslant \pi > q\)}
    {\(p > q \geqslant 0\)}
    {\(4 > p \geqslant q\)}
    {\(p \geqslant q > 0\)}
\end{problem}

\begin{answer}
B
\end{answer}

\begin{solution}
由基本不等式，\(a+\dfrac{1}{a}\geqslant2\)，所以 \(p\geqslant4\). 又因为 \(-1\leqslant t\leqslant1\)，所以 \(0\leqslant q=\arccos t\leqslant\pi\). 因此 \(p\geqslant4>\pi\geqslant q\)，从而必有 \(p>q\geqslant0\). 选 B.
\end{solution}

\section{解答题}
\begin{problem}
在直角坐标系 \(xOy\) 中，已知点 \(P(2\cos x+1, 2\cos 2x+2)\) 和点 \(Q(\cos x,-1)\)，其中 \(x \in [0,\pi]\). 若向量 \(\overrightarrow{OP}\) 与 \(\overrightarrow{OQ}\) 垂直，求 \(x\) 的值.
\end{problem}

\begin{solution}
设 \(c=\cos x\). 由 \(\overrightarrow{OP}\perp\overrightarrow{OQ}\)，得
\((2c+1)c+(2\cos2x+2)(-1)=0\). 利用 \(\cos2x=2c^2-1\)，上式化为 \(2c^2+c-2(2c^2-1)-2=0\)，即 \(c(1-2c)=0\). 因此 \(\cos x=0\) 或 \(\cos x=\dfrac{1}{2}\). 结合 \(x\in[0,\pi]\)，分别得到 \(x=\dfrac{\pi}{2}\) 或 \(x=\dfrac{\pi}{3}\). 所以 \(x=\dfrac{\pi}{3}\) 或 \(x=\dfrac{\pi}{2}\).
\end{solution}

\begin{problem}
已知实数 \(p\) 满足不等式 \(\frac{2p+1}{p+2}<0\)，试判断方程 \(z^2-2z+5-p^2=0\) 有无实根，并给出证明.
\end{problem}

\begin{solution}
由 \(\dfrac{2p+1}{p+2}<0\)，分子、分母异号，且 \(p\ne-2\)，所以 \(-2<p<-\dfrac{1}{2}\). 从而 \(p^2<4\). 视 \(z^2-2z+5-p^2=0\) 为关于 \(z\) 的一元二次方程，其判别式为 \(\Delta=(-2)^2-4(5-p^2)=4(p^2-4)<0\). 因此该方程无实根.
\end{solution}

\begin{problem}
某市 \(2003\text{ 年}\) 共有 \(1\text{ 万辆}\) 燃油型公交车. 有关部门计划于 \(2004\text{ 年}\) 投入 \(128\text{ 辆}\) 电力型公交车. 随后电力型公交车每年的投入比上一年增加 \(50\%\)，试问：

\begin{enumerate}
\item 该市在 \(2010\) 年应该投入多少辆电力型公交车？
\item 到哪一年底，电力型公交车的数量开始超过该市公交车总量的 \(\frac{1}{3}\)？
\end{enumerate}
\end{problem}

\begin{solution}
\begin{enumerate}
\item 设 \(a_n\) 表示第 \(2003+n\) 年投入的电力型公交车数量，则 \(a_1=128\)，且公比为 \(1+50\%=\dfrac{3}{2}\). 因此 \(a_n=128\left(\dfrac{3}{2}\right)^{n-1}\). \(2010\) 年对应 \(n=7\)，所以 \(a_7=128\left(\dfrac{3}{2}\right)^6=1458\)，即应投入 \(1458\text{ 辆}\).
\item 设截至第 \(2003+n\) 年底投入的电力型公交车总数为 \(S_n\). 则
\(S_n=128\left[1+\dfrac{3}{2}+\cdots+\left(\dfrac{3}{2}\right)^{n-1}\right]=256\left[\left(\dfrac{3}{2}\right)^n-1\right]\).

截至该年底，公交车总量为 \(10000+S_n\)，所以需要 \(S_n>\dfrac{1}{3}(10000+S_n)\)，即 \(S_n>5000\). 当 \(n=7\) 时，\(S_7=4118<5000\)；当 \(n=8\) 时，\(S_8=6305>5000\). 因此电力型公交车的数量到 \(2011\) 年底开始超过公交车总量的 \(\dfrac{1}{3}\).
\end{enumerate}
\end{solution}

\begin{problem}
如图，点 \(P\) 为斜三棱柱 \(ABC-A_{1}B_{1}C_{1}\) 的侧棱 \(BB_{1}\) 上一点，\(PM \perp BB_{1}\) 交 \(AA_{1}\) 于点 \(M\)，\(PN \perp BB_{1}\) 交 \(CC_{1}\) 于点 \(N\).

\begin{enumerate}
\item 求证： \(CC_{1} \perp MN\)；
\item 在任意 \(\triangle DEF\) 中有余弦定理： \(DE^2 = DF^2 + EF^2 - 2DF \cdot EF\cos \angle DFE\). 拓展到空间，类比三角形的余弦定理，写出斜三棱柱的三个侧面面积与其中两个侧面所成的二面角之间的关系，并予以证明.
\end{enumerate}

\bitmapfigure[max width=0.30\linewidth]{img/2004/shanghai_spring/q20_fig1.png}
\end{problem}

\begin{solution}
\begin{enumerate}
\item 由题意，\(PM\perp BB_1\)，\(PN\perp BB_1\)，且 \(PM\) 与 \(PN\) 相交于 \(P\)，所以 \(BB_1\perp\) 平面 \(PMN\). 又因为 \(CC_1\parallel BB_1\)，所以 \(CC_1\perp\) 平面 \(PMN\). 由于 \(MN\subset\) 平面 \(PMN\)，故 \(CC_1\perp MN\).
\item 设平面 \(BCC_1B_1\) 与平面 \(CAA_1C_1\) 所成的二面角为 \(\alpha\)，并分别记三个侧面 \(ABB_1A_1\)、\(BCC_1B_1\)、\(CAA_1C_1\) 的面积为 \(S_{AB}\)、\(S_{BC}\)、\(S_{CA}\). 由第一问知 \(CC_1\perp\) 平面 \(PMN\)，因此上述二面角的平面角为 \(\angle MNP\). 在 \(\triangle PMN\) 中，由余弦定理，
\[PM^2=PN^2+MN^2-2PN\cdot MN\cos\angle MNP.\]
两边同乘 \(CC_1^2\)，并利用 \(BB_1\parallel CC_1\)，得到
\[PM^2CC_1^2=PN^2CC_1^2+MN^2CC_1^2-2PN\cdot MN\cdot CC_1^2\cos\alpha.\]
由于 \(S_{AB}=PM\cdot BB_1=PM\cdot CC_1\)，\(S_{BC}=PN\cdot CC_1\)，\(S_{CA}=MN\cdot CC_1\)，所以
\[S_{AB}^2=S_{BC}^2+S_{CA}^2-2S_{BC}S_{CA}\cos\alpha.\]
这就是斜三棱柱三个侧面面积与其中两个侧面所成二面角之间的关系.
\end{enumerate}
\end{solution}

\begin{problem}
已知函数 \(f(x)=|x-a|\)，\(g(x)=x^2+2ax+1\)（\(a\) 为正常数），且函数 \(f(x)\) 与 \(g(x)\) 的图象在 \(y\) 轴上的截距相等.

\begin{enumerate}
\item 求 \(a\) 的值；
\item 求函数 \(f(x)+g(x)\) 的单调递增区间；
\item 若 \(n\) 为正整数，证明： \(10^{f(n)}\cdot\left(\frac45\right)^{g(n)}<4\).
\end{enumerate}
\end{problem}

\begin{solution}
\begin{enumerate}
\item 函数 \(f(x)\) 和 \(g(x)\) 在 \(y\) 轴上的截距分别为 \(f(0)=a\) 和 \(g(0)=1\). 由题意 \(a=1\).
\item 当 \(a=1\) 时，
\[
f(x)+g(x)=
\begin{cases}
x^2+x+2, & x<1,\\
x^2+3x, & x\geqslant1.
\end{cases}
\]
当 \(x<1\) 时，\(x^2+x+2=\left(x+\dfrac{1}{2}\right)^2+\dfrac{7}{4}\)，所以它在 \(\left[-\dfrac{1}{2},1\right)\) 上单调递增. 当 \(x\geqslant1\) 时，若 \(x_1<x_2\)，则 \(x_2^2+3x_2-(x_1^2+3x_1)=(x_2-x_1)(x_1+x_2+3)>0\)，所以它在 \([1,+\infty)\) 上单调递增. 两段在 \(x=1\) 处函数值相等，故 \(f(x)+g(x)\) 的单调递增区间为 \(\left[-\dfrac{1}{2},+\infty\right)\).
\item 由 \(a=1\)，且 \(n\) 为正整数，得 \(f(n)=n-1\)，\(g(n)=(n+1)^2\). 记 \(u_n=10^{f(n)}\left(\dfrac45\right)^{g(n)}=10^{n-1}\left(\dfrac45\right)^{(n+1)^2}\)，则
\[\frac{u_{n+1}}{u_n}=10\left(\frac45\right)^{2n+3}.\]
当 \(1\leqslant n\leqslant3\) 时，\(\dfrac{u_{n+1}}{u_n}\geqslant10\left(\dfrac45\right)^9=\dfrac{2621440}{1953125}>1\)，所以 \(u_1<u_2<u_3<u_4\). 当 \(n\geqslant4\) 时，\(\dfrac{u_{n+1}}{u_n}\leqslant10\left(\dfrac45\right)^{11}=\dfrac{41943040}{48828125}<1\)，所以 \(u_n\) 在 \(n\geqslant4\) 时递减. 又因为 \(\left(\dfrac45\right)^5=\dfrac{1024}{3125}<\dfrac{33}{100}\)，所以 \(u_4=10^3\left(\dfrac45\right)^{25}=10^3\left[\left(\dfrac45\right)^5\right]^5<10^3\left(\dfrac{33}{100}\right)^5=\dfrac{39135393}{10000000}<4\). 因而对任意正整数 \(n\)，都有 \(u_n\leqslant u_4<4\)，即 \(10^{f(n)}\left(\dfrac45\right)^{g(n)}<4\).
\end{enumerate}
\end{solution}

\begin{problem}
已知倾斜角为 \(45^{\circ}\) 的直线 \(l\) 过点 \(A(1,-2)\) 和点 \(B\)，\(B\) 在第一象限，\(|AB|=3\sqrt{2}\).

\begin{enumerate}
\item 求点 \(B\) 的坐标；
\item 若直线 \(l\) 与双曲线 \(C:\frac{x^2}{a^2}-y^2=1\)（\(a>0\)）相交于 \(E\)，\(F\) 两点，且线段 \(EF\) 的中点坐标为 \((4,1)\)，求 \(a\) 的值；
\item 对于平面上任一点 \(P\)，当点 \(Q\) 在线段 \(AB\) 上运动时，称 \(|PQ|\) 的最小值为 \(P\) 与线段 \(AB\) 的距离. 已知点 \(P\) 在 \(x\) 轴上运动，写出点 \(P(t,0)\) 到线段 \(AB\) 的距离 \(h\) 关于 \(t\) 的函数关系式.
\end{enumerate}
\end{problem}

\begin{solution}
\begin{enumerate}
\item 直线 \(l\) 的倾斜角为 \(45^{\circ}\)，所以斜率为 \(1\)，其方程为 \(y+2=x-1\)，即 \(y=x-3\). 设 \(\overrightarrow{AB}=(u,u)\)，由 \(|AB|=3\sqrt{2}\) 得 \(2u^2=18\)，所以 \(u=\pm3\). 因为 \(B\) 在第一象限，取 \(u=3\)，故 \(B(4,1)\).
\item 将直线 \(y=x-3\) 代入双曲线方程，得
\[\frac{x^2}{a^2}-(x-3)^2=1,\]
即 \(\left(\dfrac{1}{a^2}-1\right)x^2+6x-10=0\). 设交点 \(E\)、\(F\) 的横坐标分别为 \(x_1\)、\(x_2\). 因为线段 \(EF\) 的中点为 \((4,1)\)，所以 \(\dfrac{x_1+x_2}{2}=4\)，即 \(x_1+x_2=8\). 又因有两个交点，故 \(\dfrac{1}{a^2}-1\ne0\)，可以使用韦达定理. 于是 \(x_1+x_2=-\dfrac{6}{1/a^2-1}\)，故 \(-\dfrac{6}{1/a^2-1}=8\). 解得 \(\dfrac{1}{a^2}-1=-\dfrac{3}{4}\)，从而 \(a^2=4\). 因为 \(a>0\)，所以 \(a=2\).
\item 线段 \(AB\) 的参数方程可写为 \(Q=(1+3s,-2+3s)\)，其中 \(0\leqslant s\leqslant1\). 对于 \(P(t,0)\)，有 \(\overrightarrow{AQ}=(3s,3s)\)，\(\overrightarrow{AP}=(t-1,2)\). 当 \(\overrightarrow{AP}\cdot\overrightarrow{AB}\) 的投影落在线段 \(AB\) 内时，\(s=\dfrac{\overrightarrow{AP}\cdot\overrightarrow{AB}}{|AB|^2}=\dfrac{3(t-1)+6}{18}=\dfrac{t+1}{6}\). 因此当 \(-1\leqslant t\leqslant5\) 时，垂足在线段 \(AB\) 上，点到直线 \(AB\) 的距离为 \(\dfrac{|t-3|}{\sqrt{2}}\). 当 \(t<-1\) 时，最近点为端点 \(A\)，距离为 \(\sqrt{(t-1)^2+4}\)；当 \(t>5\) 时，最近点为端点 \(B\)，距离为 \(\sqrt{(t-4)^2+1}\). 所以
\[
h(t)=
\begin{cases}
\sqrt{(t-1)^2+4}, & t<-1,\\
\dfrac{|t-3|}{\sqrt{2}}, & -1\leqslant t\leqslant5,\\
\sqrt{(t-4)^2+1}, & t>5.
\end{cases}
\]
\end{enumerate}
\end{solution}
